% 本文件由 LaTeX 排版 Agent 根据有效 translation.json 自动生成。
% author=Hippolytus; work_id=0523; title=论圣主显

\chapter{论圣主显}

% structure: p0001 source=h1 target=omitted reason=page-title-already-rendered-by-parent

1. 我们的天主和救主的一切化工，都是美好的，甚至极为美好——凡眼睛所见、心智所悟、理智所解释、手所触摸、智慧所领悟、人性所理解的一切。因为有什么比天穹的循环更美的呢？有什么比地面的形态更绚丽的呢？有什么比太阳车的运行更快的呢？有什么比月轮更优雅的呢？有什么比星群的镶嵌更奇妙的呢？有什么比应时的风更能带来供应呢？有什么比日光更无瑕的明镜呢？有什么受造物比人更高贵呢？因此，我们的天主和救主的一切化工，都是极为美好的。再者，还有什么恩赐比水元素更为必要呢？因为万物都由水洗涤、滋养、洁净和润泽。水承载大地，水产生露水，水使葡萄树欢畅；水使谷穗成熟，水使葡萄串成熟，水使橄榄柔润，水使椰枣甜美，水使玫瑰鲜红并装饰紫罗兰，水使百合绽放它灿烂的杯盏。我何必多说呢？没有水元素，现今的秩序中无一物能存立。水元素是如此必要；因为其他元素都安置在最高穹苍之下，而水的本性却也在诸天之上获得了一席之地。先知自己就是见证，当他呼喊说：\BibleQuote{\BibleRef{咏 148:4} 天上的诸天，请赞美上主，天上的水，也请赞美上主。}\par

2. 但这还不是证明水之尊荣的唯一事物。还有比这一切更可敬的——即万有的创造者基督，曾如雨降下，被称为泉源，如河流般涌流，并在约旦河受圣洗。因为你刚听到耶稣如何来到若翰前，并由他在约旦河受圣洗。哦，无比奇妙的事！那使天主之城欢乐的无限江河，竟被浸在少许水中！那赐予众人生命、没有穷尽的无限泉源，竟被短暂贫乏的水所覆盖！那无所不在、无有不在——天使不能理解、人不能看见的那一位——按祂自己的美意来到圣洗前。亲爱的，当你听到这些事时，不要以为这是按字面说的，而要接受它们是以预象呈现的。因此，主在那隐秘的、出于仁慈屈尊就人的行为中，并未被水元素忽视。\BibleQuote{因为水一见祂，就畏惧。}它们几乎要离开原位，冲破界限。因此，先知在许多世代前就抱有这种看法，发问说：\BibleQuote{海，你为什么奔逃？约旦，你为什么倒流？}它们回答说：我们看见万有的创造者取了\InnerQuote{仆人的形状}，因不明救恩\OriginalTerm{救恩计划}{\GreekText{οἰκονομία}}的奥秘，就恐惧战栗。\par

3. 但我们知道这救恩计划的人，敬拜祂的仁慈，因为祂来是为拯救世界，而非审判世界。因此，先前不知这奥秘的主的前驱若翰，得知祂确实是主，就呼喊并对那些来受他圣洗的人说：\BibleQuote{\BibleRef{玛 3:7} 毒蛇的种类}，你们为什么这样注视我？\BibleQuote{\BibleRef{若 1:20} 我不是基督}；我是仆人，不是主人；我是臣民，不是君王；我是羊，不是牧人；我是人，不是天主。我的出生解除了我母亲的不孕；我并未使童贞不孕。我来自下界；我并未从上降下。我封住了父亲的口；我并未展开天主的恩宠。母亲认识我，我并未由星宿宣告。我卑微且微小；但\BibleQuote{\BibleRef{若 1:30} 在我以后，有一位在我以前}——按时间在我以后，但因祂不可接近、不可言喻的天主性之光，祂在我以前。\BibleQuote{\BibleRef{玛 3:11} 有一位比我更强的在我以后来，我连提祂的鞋也不配；祂要以圣神和火洗你们。}我受权威管辖，祂自身拥有权威。我被罪恶捆绑，祂是罪恶的解除者。我执行法律，祂带来恩宠。我像奴隶一样教导，祂像主人一样审判。我以大地为床，祂拥有天堂。我以悔改的圣洗施洗，祂赐予义子的恩赐：\Quote{祂要以圣神和火洗你们。}你们为什么注意我？我不是基督。\par

4. 当若翰向群众说这些话时，人们热切期待用肉眼看见某种奇观，魔鬼因若翰这样的见证而惊愕，看哪，主出现了：朴素、孤单、不加掩饰、没有随从，穿着人的身体作为外衣，隐藏着天主性的尊荣，以便躲开毒龙的陷阱。祂不仅以没有王者仪仗的主的身份来到若翰前；甚至像普通人和有罪的人一样，低头接受若翰的圣洗。因此若翰看见祂如此谦卑自己，对这件事惊异万分，开始阻止祂，说，正如你刚听到的：\BibleQuote{\BibleRef{玛 3:14} 我本受你的圣洗，而你反而到我这里来吗？}主啊，你在做什么？你所教导的不合常规。我曾宣讲过关于你的（一个道理），而你却行了另一件事；魔鬼听到了一个，却察觉到另一个。请用天主性的火洗我吧；你为什么等待水呢？请用圣神光照我；你为什么留意受造物呢？请为我这位圣洗者付圣洗，好使你的卓越性被知晓。主啊，我以悔改的圣洗施洗，且除非他们先完全告解自己的罪，我不能为他们施圣洗。那么，如果我为你施圣洗，你有什么可告解的呢？你是罪孽的解除者，还要受悔改的圣洗么？虽然我敢为你施圣洗，约旦河却不敢靠近你。\BibleQuote{\BibleRef{玛 3:14} 我本受你的圣洗，而你反而到我这里来吗？}\par

5. 主又对他说了什么？\BibleQuote{你暂且这样容许吧，因为我们应当这样履行全部的义。}\BibleRef{玛 3:15} \BibleQuote{你暂且这样容许吧，}\BibleRef{玛 3:15} \Person{若翰}；你并不比我明智。你以人的方式看，我以天主的方式预知。我应当先做这事，并这样教导。我没有任何不适当的行为，因为我被赋予了尊荣。\Person{若翰}，你惊奇我不是以我的尊严而来的吗？君王的紫袍不适合平民，但军人的荣耀适合君王：我是来见一位王子，而不是一位朋友吗？\BibleQuote{你暂且这样容许吧，因为我们应当这样履行全部的义：}\BibleRef{玛 3:15} 我是法律的完成者；我寻求不留下任何未完成之处，以便在我之后，\Person{保禄}可以宣告：\BibleQuote{基督是法律的终结，使凡信者获得正义。}\BibleRef{罗 10:4} \BibleQuote{你暂且这样容许吧，因为我们应当这样履行全部的义。}\BibleRef{玛 3:15} \Person{若翰}，为我洗吧，好使没有人轻视圣洗。由你——仆人——受洗，好使君王或权贵中没有人不屑于由一位贫穷司铎的手受洗。让我下到\Place{约旦河}，好使他们能听到我父的见证，并承认圣子的权能。\BibleQuote{你暂且这样容许吧，因为我们应当这样履行全部的义。}\BibleRef{玛 3:15} 于是\Person{若翰}终于允许了他。\BibleQuote{耶稣受了洗，立刻从水里上来；天为他开了；他看见天主的圣神仿佛鸽子降下，落在他身上。又有声音从天上说：这是我的爱子，我所喜悦的。}\BibleRef{玛 3:16-17}\par

6. 亲爱的，你看见了吗？如果主听从了\Person{若翰}的劝告而拒绝受洗，我们会失去多少、多大的祝福？因为在这之前，天是关闭的；上界是难以进入的。那样我们就会降到下层，却不能升到上层。但难道主只是受了洗吗？他也更新了旧人，并再次将嗣子的权杖交给了他。因为立刻\BibleQuote{天为他开了}。可见与不可见之间发生了和好；天上的秩序充满喜乐；地上的疾病得医治；隐秘的事被显明；敌对者复归于友好。因为你已听到福音作者的话，他说\BibleQuote{天为他开了}，是由于三个奇事。因为当新郎基督受洗时，天上的洞房应当敞开它灿烂的门。同样，当圣神以鸽子的形状降下，圣父的声音传遍各处时，也当\Quote{天门的门扉被举起}。\BibleQuote{看，天为他开了，又有声音说：这是我的爱子，我所喜悦的。}\BibleRef{玛 3:16-17}\par

7. 爱者生爱，无形之光生不可接近之光。 \BibleQuote{这是我的爱子}\BibleRef{玛 3:17}，他显现于地上却未与圣父的怀抱分离，被彰显却未显现。因为显现是一回事，在此处施洗者似乎大于受洗者。为此，圣父从天上派遣圣神降在受洗者身上。因为在\Person{诺厄}的方舟中，天主对人的爱由鸽子所象征，同样，如今圣神以鸽子的形状降下，仿佛衔着橄榄的果实，落在那位被见证者身上。为了什么？为使圣父声音的可靠显明，并使许久以前的预言得以证实。那预言是什么？\BibleQuote{主的声音在大水之上，荣耀的天主打雷；主在大水之上。}\BibleRef{咏 29:3} 这是什么声音？\BibleQuote{这是我的爱子，我所喜悦的。}\BibleRef{玛 3:17} 他就是那被称为\Person{若瑟}的儿子，按神圣本体是我的独生子。 \BibleQuote{这是我的爱子}\BibleRef{玛 3:17}——他饥饿，却维持着无数人；他疲倦，却赐予疲倦者安息；他没有枕头的地方，却用手托住万有；他受难，却医治苦难；他被打，却将自由赐给世界；他被刺穿肋旁，却修复了\Person{亚当}的肋旁。\par

8. 但请现在给我你最好的注意，我恳求你，因为我愿回到生命的泉源，观看那涌出医治的泉源。不朽的圣父派遣了不朽的圣子和圣言进入世界，他来到人那里，为用水和圣神洗净他；他重生我们，使灵魂与身体不朽，将生命的呼吸（圣神）吹入我们内，并赐予我们不朽的全副武装。因此，如果人成为不朽的，他也会成为神。如果他在水与圣神中，在圣洗池的再生之后成为神，那么在死后复活后，他也被发现是基督的共同继承人。因此，我宣讲如下：\Quote{来吧，万民的所有宗族，到圣洗的不朽那里。我将生命的喜讯带给你们这些停留在无知黑暗中的人。从奴役进入自由，从暴政进入王国，从腐朽进入不朽。} 有人说，我们如何来？如何？藉着水和圣神。这就是与圣神结合的水，由此乐园得到浇灌，大地得到滋养，植物生长，动物繁衍，并且（用一句话概括）人借以重生并获得生命，基督也在这水中受洗，圣神以鸽子的形状降下。\par

9. 这就是那起初\BibleQuote{运行在水面上}的圣神，\newline{}世界藉祂而运行，万有藉祂而存在，一切生物藉祂而有生命；\newline{}祂也大能地运行在先知们身上，并降落在基督身上。\newline{}这就是那以火舌形式赐给宗徒们的圣神。\newline{}这就是达味所寻求的圣神，他曾说：\BibleQuote{天主，求祢给我再造一颗纯洁的心，在我五内注入坚强的精神。}\newline{}关于这圣神，加俾额尔也曾对童贞玛利亚说：\BibleQuote{圣神要临于你，至高者的能力要庇荫你。}\newline{}藉这圣神，伯多禄说出了那句有福的话：\BibleQuote{祢是基督，永生天主之子。}\newline{}藉这圣神，教会的磐石得以建立。\newline{}这就是那护慰者圣神，祂为了你们而被派遣，为向你们显示你们是天主之子。\par

10. 那么，人啊，来，再生于天主的义子之列。\newline{}有人问：如何作到？\newline{}如果你不再犯奸淫，不杀人，不事奉偶像；如果你不被享乐所制伏；如果你不让骄傲的情感统治你；如果你除去不洁的污秽，卸下罪恶的重担；如果你丢弃魔鬼的武器，穿上信德的胸甲，正如依撒意亚所说：\BibleQuote{你们要洗涤，寻求正义，解救受压迫的人，为孤儿伸冤，为寡妇辩护。然后来，我们彼此辩论——上主说。你们的罪虽似朱红，我将使它们白如雪；虽红如丹颜，我将使它们白如羊毛。如果你们愿意，听从我的声音，将享用地上的美物。}\newline{}亲爱的，你看见先知如何预言了圣洗的净化之能吗？\newline{}因为那怀着信德下到重生池中的人，弃绝魔鬼，与基督结合；他否认仇敌，宣认基督是天主；他卸下奴役，穿上义子的身份——他从圣洗中上来，光明如太阳，闪耀着正义的光辉，而最重要的是，他返回时已成为天主的儿子，基督的共同继承人。\newline{}愿光荣和权能归于祂，并与祂至圣、至善、赋予生命的圣神，从今世直到永远，万世万代。阿们。\par

\SourceRecord{Translated by J.H. MacMahon.；Ante-Nicene Fathers；Vol. 5.；Edited by Alexander Roberts, James Donaldson, and A. Cleveland Coxe.；Buffalo, NY: Christian Literature Publishing Co.,；1886.；<http://www.newadvent.org/fathers/0523.htm>.}
